\documentclass[output=book,minimal,biblatexbackend=bibtex]{langscibook}

\title{Notes on Open-Access Book Production}
\subtitle{A worked example}
\author{Noor Fatima}
\BackBody{A short example of the langscibook class, used by Language Science
Press for open-access linguistics monographs.}

\addbibresource{references.bib}

\begin{document}

% langscibook's `minimal` mode (chosen because it skips loading fonts this
% image does not have — Libertinus, XITS, Arimo) also redefines \maketitle to
% print only "You are using the minimal mode.", since the mode exists for
% fast draft compiles rather than a finished cover. That placeholder is fine
% mid-draft but wrong as the first thing a downloaded starter shows, so it is
% overridden back to a plain title block here.
\makeatletter
\renewcommand{\maketitle}{%
  \begingroup
  \centering
  \vspace*{3cm}
  {\LARGE\bfseries \@title\par}
  \vspace{0.5cm}
  {\large \@subtitle\par}
  \vspace{1.5cm}
  {\large \@author\par}
  \endgroup
  \clearpage
}
\makeatother
\maketitle
\frontmatter
\tableofcontents

\mainmatter

\chapter{Preparing a Manuscript}

\section{Why a Press-Specific Class}

A publisher-specific class encodes house style directly, so authors do not
have to reproduce it by hand \citep{knuth1984texbook}. \texttt{langscibook}
covers title pages, running heads and a bibliography style consistent across
every book in the series.

\section{Structure of This Example}

Table~\ref{tab:sections} outlines the two chapters in this starter and what
each demonstrates.

\begin{table}[h]
  \centering
  \begin{tabular}{ll}
    \toprule
    Chapter & Demonstrates \\
    \midrule
    1 & Sectioning and a table \\
    2 & Citation and bibliography \\
    \bottomrule
  \end{tabular}
  \caption{What each chapter in this starter covers.}
  \label{tab:sections}
\end{table}

\chapter{Citing Sources}

\section{A Resolved Citation}

The reference below is resolved by \texttt{biblatex} in bibtex-backend mode,
which keeps the compile loop the same as a plain \texttt{book} document
\citep{knuth1984texbook}.

\section{Closing Remarks}

The two chapters above are enough to show a working title page, table of
contents, sectioning, a table and a bibliography under
\texttt{langscibook}'s \texttt{minimal} mode, which skips cover generation and
does not load font files.

\printbibliography

\end{document}
